\subsection{Control law}
Like in \cite{Cas:18} the GSTA algorithm is given by
\begin{equation}
    \begin{split}
        u &= -k_1 \phi_1(\sigma) + z \\ 
        \dot{z} &= -k_2\phi_2(\sigma) \\
        \phi_1(\sigma) &= \lceil \sigma \rfloor^{1/2} + \beta \sigma \\
        \phi_2(\sigma) &= \frac{1}{2} \lceil \sigma \rfloor^0 + \frac{3}{2}\beta\lceil \sigma \rfloor^{1/2} + \beta^2\sigma  
    \end{split}
\end{equation}

\subsubsection{Adaptive gain law}
Using same adaptive law as \cite{Sht:10}
\begin{equation}
    \begin{split}
    \dot{k}_1 &= 
    \begin{cases}
        &\omega_1 \sqrt{\frac{\gamma_1}{2}}, \text{ if } \sigma \neq 0, \\
        &0, \qquad\ \ \text{ if } \sigma = 0
    \end{cases} \\
    k_2 &= 2\epsilon k_1 + \lambda + 4\epsilon^2
    \end{split}
    \label{eq:borlaug-adaptive-law}
\end{equation}

\subsubsection{Closed-loop system}
    \begin{equation}
        \dot{\sigma} = -k_1\gamma(\sigma, t)\phi_1(\sigma) + \varphi_1(\sigma, t) + \gamma(\sigma, t)\left(z + \frac{\varphi_2(\sigma, t)}{\gamma(\sigma, t)}\right)
    \end{equation}
    the system can be broken down
    \begin{equation}
        \begin{split}
        \dot{\sigma}_1 &= \gamma(\sigma, t)\left(-k_1\phi_1(\sigma_1) + \frac{\varphi_1(\sigma_1,t)}{\gamma(\sigma, t)} + \sigma_2\right)\\
        \dot{\sigma}_2 &= -k_2\phi_2(\sigma_1)+ \frac{d}{dt}\left(\frac{\varphi_2(\sigma, t)}{\gamma(\sigma, t)}\right)
        \end{split}
        \label{eq:borlaug-closed-loop}
    \end{equation}
    where $\sigma_1 = \sigma$ and $\sigma_2 = z + \frac{\varphi_2(\sigma, t)}{\gamma(\sigma, t)}$

\begin{theorem}[Borlaug et al. \cite{Bor:22}]
    Suppose that $\gamma(\sigma_1, t)$ and $\varphi(\sigma_1, t)$ in system \eqref{eq:castillo-general-case} satisfy Assumptions 1-4, from \cite{Cas:18}. Then, the closed-loop dynamics \eqref{eq:borlaug-closed-loop} are globally finite-time stable, such that the states $\sigma_1$ and $\sigma_2$ converge to zero and z converges to $-\varphi(0, t)/\gamma(0,t)$, globally and in finite time, ig the gains $k_1$ and $k_2$ are designed as expressed in \eqref{eq:borlaug-adaptive-law}, $\beta > 0, \lambda > 0, \omega_1 > 0, \gamma_1 > 0,$ and $\epsilon= \frac{\omega_2}{2\omega_1}\sqrt{\frac{\gamma_2}{\gamma_1}}$, where $\omega_2 > 0$ and $\gamma_2 > 0$.
\end{theorem}
\begin{rmk}
    \begin{itemize}
        \item The theorem and proof can be extended to $n$ dimensions by separating the problem into $n$ one-dimensional cases
        \item GSTA is considered proven by Castillo et al. \cite{Cas:18}
        \item The rest of the proof is similar to that of Shtessel et al. \cite{Sht:10}
    \end{itemize}
\end{rmk}

\subsection{Real system}
Equations of motion
    \begin{equation}
    \begin{split}
        \dot{\xi} &= J(\eta_2)\zeta = 
    \begin{bmatrix}
        R_{Ib}^\top(\eta_2) && 0_{3 \times 3} && 0_{3 \times (n-1)} \\
        0_{3 \times 3} && J^{-1}_{k,o}(\eta_2) && 0_{3 \times (n-1)} \\
        0_{(n-1) \times 3} && 0_{(n-1) \times 3} && I_{(n-1) \times (n-1)}
    \end{bmatrix}\zeta \\
    M&(q)\dot{\zeta} + C(q, \zeta)\zeta + D(q, \zeta)\zeta + g(q, R_{Ib}) = \tau(q)
    \end{split}
\end{equation}
where $\xi = \left[\eta_1^\top, \eta_2^\top, q^\top\right]^\top$, here $\eta_1 = \left[x, y, z\right]^\top \in \R^{6+(n-1)}$ is the position, $\eta_2 = \left[\phi, \theta, \psi\right]^\top \in \R^3$ are Euler angles and $q \in \R^{n-1}$ are the joint angles. The velocity vector $\zeta = \left[\nu^\top, \omega^\top, \dot{q}^\top\right]^\top \in \R^{6+(n-1)}$ comprises body-fixed linear and angular vehicle velocities $\nu$ and $\omega$, and joint velocities $\dot{q}$.


\subsubsection{Tracking control law using AGSTA}
We choose sliding manifold
\begin{equation}
    \sigma = \zeta - \zeta_r \in R^{6+(n-1)}
\end{equation}
where $\zeta$ are velocities of the body and $\zeta_r$ is a velocity reference trajectory. The control input is given by
\begin{equation}
    \tau(q) = u_{\text{AGSTA}} \in R^{6+(n-1)}
\end{equation}

The derivative of the sliding variable is 
\begin{equation}
    \dot{\sigma} = M^{-1}(\cdot)\left[-C(\cdot)(\zeta)-D(\cdot)(\zeta) - g(\cdot) + \tau(\cdot)\right] - \dot{\zeta}_r
\end{equation}
using $\zeta = \zeta_r + \sigma$
\begin{equation}
    \dot{\sigma} = M^{-1}(\cdot)\left[-C(\cdot)(\zeta_r + \sigma)-D(\cdot)(\zeta_r + \sigma) - g(\cdot) + \tau(\cdot)\right] - \dot{\zeta}_r
\end{equation}
\begin{rmk}
    If time: Can Kristin explain why $(\cdot)$ was used to replace the arguments? 
\end{rmk}

The model is then divided into it's respective degrees of freedom and $\tau(q) = u_{AGSTA}$ is applied.
% TODO: write down the AIAUV controller for this

\subsection{Experiments}
Tracking performance \cite{Bor:22}
% \begin{figure}[h]
%     \includegraphics[width=\textwidth]{figures/borlaug-tracking.png}
% \end{figure}
Gain evolution \cite{Bor:22}
% \begin{figure}[h]
%     \includegraphics[width=\textwidth]{figures/borlaug-gains.png}
% \end{figure}

\subsection{Comparison}
Performance \cite{Bor:22}
% \begin{figure}[h]
%     \includegraphics[width=\textwidth]{figures/borlaug-table-1.png}
% \end{figure}

Actuation \cite{Bor:22}
% \begin{figure}[h]
%     \includegraphics[width=\textwidth]{figures/borlaug-table-2.png}
% \end{figure}



\subsection{Key takeaways}
Features
\begin{itemize}
    \item Deals well with matched uncertainties
    \item Works for uncertain input gain
    \item Practically proven control method with desirable stability guarantees
    \item AGSTA demonstrates better tracking performance with GSTA with constant gains
\end{itemize}
Future work
\begin{itemize}
    \item Coupled MIMO formulation
    \item Adaptive law with fewer parameters
    \item Adaptive law that enables decreasing controller gain
\end{itemize}
